\documentclass[11pt]{tex/lawoutline}
\usepackage{amssymb}

\course{Criminal Law}
\student{Prof: I. Rabb}
\semester{Fall 2026}

\begin{document}

\maketitle
\tableofcontents

\chapter{Background}

\chapter{\textit{Actus Reus}}

\section{Voluntary Acts}

\section{Acts of Omission}

\section{Police Duty to Intervene}

\chapter{\textit{Mens Rea}}

\section{Basic Concepts}

\section{Strict Liability}

\section{Negligence, Statutory Silence}

\chapter{Causation and Attendant Circumstances}

\section{Causation}

\subsection{Foreseeability and Coincidence}

\definition{Proximate Cause}{Historically used to separate results that an actor would be held responsible for from results an actor would not be held responsible for. The line of demarcation is flexible, and attempts at uniformity have failed, but that doesn't mean general guidelines and analysis approaches cannot be constructed.}

\rulebox{\textit{sine qua non} test}{An essential condition, action, or ingredient without which something else is impossible or cannot exist. (But-for test)}

\casebox
    {\textit{People v. Acosta} (1991)}
    {\textit{People v. Acosta}, 232 Cal. App. 3d 1375, 284 Cal. Rptr. 117 (Cal. Ct. App. 1991).}
    {
        \textbf{Facts} 
        \begin{itemize}
            \item Acosta was apprehended by police in a stolen vehicle
            \item Acosta led police on a high-speed chase through Orange County (48 miles)
            \item Acosta drove "egregiously:" running stop signs and red lights, driving on the wrong side of the street, etc.
            \item Police helicopters helped in tracking Acosta during the chase
            \item Two police helicopters collided, killing three people (police officers)
            \item Turner (expert) says that accident happened because one pilot violated FAA regulations
            \item Acosta argues that $\nexists$ sufficient evidence proving he proximately caused the death of his victims, because a crash between airborne vehicles was not a foreseeable result of his conduct
        \end{itemize}
        
        \textbf{Procedural History}
        \begin{itemize}
            \item Trial court jury convicted Acosta of three counts of 2nd degree murder and one count of unlawfully driving or taking a vehicle without consent
            \item Acosta appealed his conviction of 2nd degree murder in California's Court of Appeals (4th district)
        \end{itemize}
        
        \textbf{Relevant law}
        \begin{itemize}
            \item Proximate cause of 2nd degree murder
        \end{itemize}
        
        \textbf{Question}
        \begin{itemize}
            \item Was Acosta's actions the "actual cause" of the victim's injury? 
            \item Equivalently: But for Acosta's actions, would the injury have occurred? 
        \end{itemize}
        
        \textbf{Holding}
        \begin{itemize}
            \item Unless an act is an actual cause of the injury, it will not be considered a proximate cause. 
            \item $\exists$ proximate cause here
        \end{itemize}
        
        \textbf{Discussion (J. Wallin)}
        \begin{itemize}
            \item Virtually impossible to have a logical verbal formula that produces uniform results
            \item Standard should be to exclude extraordinary results and determination of facts to be through "the common sense of the common man as to common things"
            \item Definition of foreseeability in California: not involving $\Delta$'s state of mind, but rather the objective conditions present when he acts
            \item But for Acosta fleeing the police, the helicopters wouldn't have been in the position for a crash
            \item Result of helicopter crash is not highly extraordinary
            \item Given emotional dynamics of pursuit, $\exists$ "appreciable probability" that pursuer could act negligently or recklessly in order to catch the $\Delta$ 
            \item Mens rea: the court found insufficient proof of malice for a murder conviction, as in Acosta didn't consciously disregard the risk to the helicopter pilots
        \end{itemize}
        
        \textbf{Dissent (J. Crosby)}
        \begin{itemize}
            \item The law does not assign blame to an otherwise blameworthy actor when neither the intervening negligent conduct nor the risk of harm was foreseeable
            \item Occupants of helicopter were not within the range of apprehension of a fleeing criminal on the ground
            \item This \textit{was} a highly extraordinary result
        \end{itemize}
    }

\casebox
    {\textit{People v. Arzon} (1978)}
    {\textit{People v. Arzon}, 92 Misc. 2d 739, 401 N.Y.S.2d 156 (Sup. Ct. N.Y. County 1978).}
    {
        \textbf{Facts}
        \begin{itemize}
            \item $\Delta$ allegedly set fire to a couch on the fifth floor of an abandoned building
            \item Firefighters responded but withdrew after failing to control the fire
            \item During evacuation a separate fire caused smoke that made evacuation dangerous
            \item Firefighter Celic was injured during the evacuation and later died
            \item The second fire was also arson, but not sufficient evidence that connected Arzon to it
            \item Arzon charged with depraved-indifference murder and felony murder
        \end{itemize}

        \textbf{Procedural History}
        \begin{itemize}
            \item $\Delta$ indicted for two counts of 2nd degree murder and one count of 3rd degree arson.
            \item $\Delta$ moved to dismiss the two murder counts, arguing grand-jury evidence was insufficient and that prosecution had not established a causal relationship with his actions and Celic's death
        \end{itemize}

        \textbf{Question}
        \begin{itemize}
            \item Was there sufficient relationship between the $\Delta$'s actions and Celic's death?
            \item Court: Yes
        \end{itemize}

        \textbf{Holding}
        \begin{itemize}
            \item The $\Delta$'s conduct doesn't need to be the sole and exclusive factor in the victim's death
            \item The fire set by the $\Delta$ was an indispensable link in the chain of events that caused Celic's death
        \end{itemize}

        \textbf{Discussion}
        \begin{itemize}
            \item \textit{People v. Kibbe}: $\Delta$s' conduct was a sufficiently direct cause of victim's ensuing death
            \item Similar to \textit{Kibbe} and how the $\Delta$s were still charged with murder even though the victim died in a separate incident
            \item Foreseeable that the firemen would respond to the situation ($\Delta$'s fire) thus exposing them to life-threatening danger
            \item $\Delta$'s act put Celic in a particularly vulnerable position
        \end{itemize}
    }

    \casebox
        {\textit{People v. Warner-Lambert Co.} (1980)}
        {\textit{People v. Warner-Lambert Co.}, 414 N.E.2d 660 (N.Y. 1980)}
        {
            \textbf{Facts}
            \begin{itemize}
                \item Warner-Lambert indicted for 2nd degree manslaughter and criminally negligent homicide
                \item Several employees killed during an explosion at a factory 
                \item $\exists$ evidence that corporation used magnesium stearate, which is potentially explosive, in its manufacturing process
                \item $\Delta$ knew about explosion hazard
                \item Speculation on what caused the explosion
            \end{itemize}

            \textbf{Holding}
            \begin{itemize}
                \item Court held that there was insufficient evidence to establish foreseeability of the immediate cause of the explosion
                \item Dismissed the indictment
            \end{itemize}

            \textbf{Discussion}
            \begin{itemize}
                \item Reject prosecution's broad application of proximate cause
                \item Prosecution's argument is that exposure of risk to death $\rightarrow$ death should be criminally liable
                \item $\Delta$'s actions must be a \textit{sufficiently direct cause} of the ensuing death
            \end{itemize}
        }

\note
{Statutory Standards}
{
    \begin{itemize}
        \item Most state codes have no explicit rules for determining causation
        \item Courts are left to decide based on evolving common-law principles
    \end{itemize}
}

\note
{The Specific Causal Mechanism}
{
    \begin{itemize}
        \item Recently: NY courts limited the specific-causal-mechanism requirement to only contexts that involve a commercial or manufacturing process (not to fires in residential settings, etc.)
        \question{Why not?}{If the specific-causal-mechanism requirement is sound, shouldn't it apply in residential situations as well?}
    \end{itemize}
}

\note
{Criticisms of Foreseeability Test}
{
    \begin{itemize}
        \item Concept of foreseeability is inherently manipulable because of subjectivity in description of events (Prof. Moore)
        \item Alternatively, depend less on the degree of foreseeability or risk and more on the relationship between $\Delta$'s conduct and risk arising from it (J. Posner)
    \end{itemize}
}

\rulebox
{The Vulnerable Victim}
{
    \begin{itemize}
        \item When a victim with underlying medical conditions died of fright, $\Delta$ was found guilty of murder, and the court said that liability was not limited to the deaths that were foreseeable
        \item But, if the victim dies \textit{after} an assault, $\Delta$ is not found guilty
        \question{Fairness}{Is this fair?}
    \end{itemize}
}


\end{document}